\documentclass[11pt,letterpaper]{article}
\usepackage[top=0.42in,bottom=0.45in,left=0.45in,right=0.45in]{geometry}
\usepackage[T1]{fontenc}
\usepackage[utf8]{inputenc}
\usepackage{mathptmx}
\usepackage{microtype}
\usepackage{enumitem}
\usepackage{titlesec}
\usepackage{needspace}
\usepackage[hidelinks]{hyperref}
\hypersetup{pdftitle={Pratheek Eranki Resume},pdfauthor={Pratheek Eranki}}
\pagestyle{empty}
\raggedright
\setlength{\parindent}{0pt}
\setlength{\parskip}{1.5pt}
\setlength{\emergencystretch}{2em}
\linespread{0.98}
\titleformat{\section}{\needspace{3\baselineskip}\fontsize{12}{14}\selectfont\bfseries}{}{0pt}{}[\vspace{-1pt}\titlerule]
\titlespacing*{\section}{0pt}{4pt}{4pt}
\setlist[itemize]{leftmargin=11pt,itemsep=1pt,topsep=2pt,parsep=0pt,partopsep=0pt,label=\textbullet}
\newcommand{\entry}[3]{\needspace{11\baselineskip}\noindent\textbf{#1}\hfill #3\\[-1pt]\textit{#2}\par}
\begin{document}
\fontsize{10.5}{12}\selectfont
\begin{center}
{\fontsize{29}{32}\selectfont Pratheek Eranki}\par
\vspace{4pt}

{\small{} St. Charles, MO \textbar{} \href{mailto:pratheek.eranki@gmail.com}{pratheek.eranki@gmail.com} \textbar{} \href{tel:+19019028689}{+1 901-902-8689} \textbar{} \href{https://www.linkedin.com/in/pratheekeranki64/}{LinkedIn} \textbar{} \href{https://github.com/PratheekEranki}{GitHub} \textbar{} \href{https://pratheekeranki.github.io/portfolio/}{Portfolio}}\par\end{center}\vspace{-9pt}
\section*{Summary}
Software engineer with experience in enterprise application delivery, sleep EEG and fMRI research data preparation, and legislative data analysis. Builds full-stack applications and AI/data projects using Python, Java, SQL, React, and machine learning tools. Co-author of the Rhamba research publication.
\section*{Education}
\textbf{University of Memphis}\hfill{\small{}Aug 2024 -- May 2026}\\
M.S. in Data Science\hfill{\small{} GPA 3.87 / 4.0}\par\vspace{2pt}
\section*{Skills}
\textbf{Languages:} Python, SQL, Java, JavaScript, TypeScript, R, Scala\par
\textbf{Deep learning \& architectures:} PyTorch, TensorFlow, Transformers / self-attention, Mamba / state space models, Self-supervised learning, Masked autoencoders, Model fine-tuning, CBraMod\par
\textbf{LLMs \& retrieval:} Hugging Face Transformers, Sentence Transformers, RAG, FAISS, LangChain, NVIDIA NIM, LLM APIs (OpenAI, Anthropic, Gemini), Prompt engineering\par
\textbf{Machine learning \& evaluation:} scikit-learn, XGBoost, LightGBM, Time-series forecasting, Cross-validation, SHAP, BLEU / ROUGE / F1 evaluation\par
\textbf{Research \& ML compute:} Resting-state fMRI, EEG sleep staging, MNE-Python, CUDA, Multi-GPU training, HPC clusters\par
\textbf{Data \& visualization:} pandas, NumPy, SciPy, ETL pipelines, Apache Spark, Tableau, Plotly, Matplotlib, Seaborn, Excel\par
\textbf{Backend \& APIs:} Spring Boot, Spring Security, FastAPI, REST APIs, JWT authentication, WebSockets, Apache Kafka\par
\textbf{Frontend \& applications:} React, HTML, CSS, Tailwind CSS, Streamlit, Gradio\par
\textbf{Databases \& persistence:} PostgreSQL, MySQL, MongoDB, Redis, SQLAlchemy, Spring Data JPA, Relational database design, Query optimization\par
\textbf{MLOps, delivery \& practices:} MLflow, Evidently, AWS, Docker, Docker Compose, Git, GitHub Actions / CI/CD, JUnit, pytest, Agile / Scrum, Technical documentation\par
\section*{Experience}
\entry{St. Jude Children's Research Hospital}{Software Engineer}{September 2025 -- May 2026}
Contributed to software and data workflows supporting sleep EEG and fMRI research in a clinical research environment.
\begin{itemize}
\item Prepared and validated sleep EEG and fMRI data for downstream research analysis, resolving data-quality and formatting issues while maintaining traceable processing steps.
\item Used PyTorch and TensorFlow in research workflows involving Transformers, self-attention, and Mamba architectures; supported model preparation, training, and fine-tuning experiments.
\item Worked with CBraMod-based EEG sleep-staging workflows and self-supervised fMRI experiments, connecting data preparation with model inputs and downstream evaluation.
\item Collaborated with researchers to translate study requirements into repeatable processing and analysis workflows, and developed utilities to reduce repetitive preparation work.
\item Documented processing procedures, validation checks, experiment configurations, and technical decisions so research workflows could be reproduced, reviewed, and maintained.
\end{itemize}\vspace{2pt}
\entry{University of Memphis}{Graduate Student Research Assistant}{March 2025 -- August 2025}
Conducted data-focused research using legislative records obtained through LegiScan.
\begin{itemize}
\item Built reproducible workflows to collect legislative records through the LegiScan API, clean source data, and organize structured SQL datasets for research analysis.
\item Standardized bill identifiers, jurisdictions, legislative status, sponsors, and subject classifications to support consistent comparisons across legislative records.
\item Used SQL and Excel to clean, deduplicate, and align legislative records with demographic data by state and year for comparative analysis.
\item Performed exploratory analysis of legislative activity and prepared Tableau and Plotly views to communicate trends and structured findings for faculty review.
\item Documented data definitions, transformation logic, API collection steps, and known limitations so research workflows and datasets could be reproduced and extended.
\end{itemize}\vspace{2pt}
\entry{Cognizant}{Software Engineer}{October 2021 -- April 2024}
Contributed to enterprise software delivery in client-facing engineering environments.
\begin{itemize}
\item Developed and maintained Java Spring Boot application components and REST APIs within established development, testing, and release processes.
\item Built and maintained React and TypeScript interface components for data display, form validation, and application workflows, integrating them with backend services.
\item Worked with MySQL queries and indexing to support application data access and reporting; investigated software defects, documented root causes, and coordinated fixes with engineering and business stakeholders.
\item Supported JUnit testing, Docker-based environments, AWS deployment workflows, CI/CD, release validation, and production issue resolution across iterative software deliveries.
\item Participated in Agile delivery and code reviews, and created technical documentation and handoff materials explaining application behavior, recurring issues, and maintenance procedures.
\end{itemize}\vspace{2pt}
\section*{Projects}
\needspace{5\baselineskip}\textbf{FieldVerify}\hfill{\small{} \href{https://github.com/PratheekEranki/fieldverify}{Code} \textbar\ \href{https://pratheekeranki.github.io/portfolio/\#project/fieldverify}{Case study}}\\
{\small{}FastAPI / React / Android / Identity SDKs}\par
\begin{itemize}
\item Built the workflow around separate authenticity and matching signals, manual fallback, SDK integration, and a chained audit log.
\item Uses a FastAPI backend with React and Android interfaces.
\item Separates cryptographic credential checks from record-match scoring.
\item Records outcomes in a SHA-256-chained audit trail and includes tests for matching, audit behavior, and verification paths.
\end{itemize}\vspace{3pt}
\needspace{5\baselineskip}\textbf{OpsFlow}\hfill{\small{} \href{https://github.com/PratheekEranki/opsflow-saas}{Code} \textbar\ \href{https://pratheekeranki.github.io/portfolio/\#project/opsflow}{Case study}}\\
{\small{}Spring Boot / React / PostgreSQL / WebSockets}\par
\begin{itemize}
\item Implemented team and ticket workflows across Spring Boot services and the React interface, including role-aware access and real-time updates.
\item Scopes organizations, users, projects, and tickets through authenticated APIs.
\item Uses WebSockets for board updates and Kafka for activity events.
\item Combines PostgreSQL persistence, Redis caching, and container-based development.
\end{itemize}\vspace{3pt}
\needspace{5\baselineskip}\textbf{FixBot}\hfill{\small{} \href{https://github.com/PratheekEranki/fixbot-device-agent}{Code} \textbar\ \href{https://pratheekeranki.github.io/portfolio/\#project/fixbot}{Case study}}\\
{\small{}Python / NVIDIA NIM / FAISS / FastAPI / Gradio}\par
\begin{itemize}
\item Implemented a custom iterative tool-calling loop, structured lookup tools, session memory, FastAPI endpoints, and a Gradio interface showing tool activity.
\item Dispatches three tools for symptoms, pricing, and FAISS-backed knowledge search.
\item Feeds tool results back into an OpenAI-compatible NVIDIA NIM model until a response is ready.
\item Keeps a sliding window of conversation history and exposes the API separately from the interface.
\end{itemize}\vspace{3pt}
\section*{Publications}
Pankaj Pandey, Ruthwik Reddy Doodipala, Pratheek Eranki, Carolina Torres-Rojas, Manob Jyoti Saikia, Ranganatha Sitaram. \textit{Rhamba: Region-Aware Hybrid Attention-Mamba Framework for Self-Supervised Learning in Resting-State fMRI}. Accepted for publication in Computers in Biology and Medicine, as reported by the arXiv record. \href{https://arxiv.org/abs/2605.01240}{arXiv:2605.01240}.
\section*{Certifications}
\begin{itemize}
\item AWS Certified Cloud Practitioner \textemdash{} Amazon Web Services \hfill {\small{}Nov 2023 \textperiodcentered{} Expires Nov 2026}
\item Microsoft Certified: Azure Fundamentals \textemdash{} Microsoft \hfill {\small{}Apr 2023}
\item AI Fluency: Framework \& Foundations \textemdash{} Anthropic \hfill {\small{}Jun 2026}
\item Crash Course on Python \textemdash{} Coursera \hfill {\small{}Jun 2020}
\end{itemize}
{\small{} \href{https://www.linkedin.com/in/pratheekeranki64/details/certifications/}{Credentials on LinkedIn}}
\end{document}
